% GENERATED FILE. Edit canonical lessons, then run npm run generate:books.
% canonical-source-hash: fnv1a64:533561de577f6fcd
% canonical-lessons: SA-C221-karpasah, SA-C221-kauseyam, SA-C221-carma, SA-C221-granthih, SA-C221-pithika

\chapter{Cotton, Silk, Leather, Knot, Small stool}
\label{ch:sa-cotton-silk-leather-knot-small-stool}
\hlchaptermodality{221}

\begin{chapteropening}
\textbf{By the end of this chapter:} \emph{I can say cotton, silk, leather, a knot and a small stool.}

Complete the last lesson of chapter 221: I can say cotton, silk, leather, a knot and a small stool.
\end{chapteropening}

\section[\texorpdfstring{\sk{कार्पासः} (kārpāsaḥ)}{कार्पासः (kārpāsaḥ)}]{\sk{कार्पासः} (kārpāsaḥ) — cotton}

\label{lesson:SA-C221-karpasah}



\begin{quote}
Before the new one: say the Sanskrit for a staff, then the Sanskrit for an archer's bow.
\end{quote}

\subsection*{You'll want to know: \sk{कार्पासः}}
\textbf{\sk{कार्पासः}} — \emph{kārpāsaḥ} — \textquotedblleft{}cotton\textquotedblright{}.

\begin{grammarlens}[title={What you've built}]
One more word for this chapter.
\end{grammarlens}

\subsection*{Your turn}
\begin{itemize}
\raggedright
  \item \emph{Say it:} \emph{kārpāsaḥ}
  \item \emph{Say it:} \emph{kārpāsaḥ}, once more
  \item \emph{Say it:} say \emph{kārpāsaḥ}
  \item \emph{Recall:} say the Sanskrit for a staff, then the Sanskrit for an archer's bow, then say \emph{kārpāsaḥ} again
\end{itemize}

\begin{tcolorbox}[breakable,colback=teal!4,colframe=teal!35!black,title={Before you move on}]
What does \sk{कार्पासः} mean? (\textquotedblleft{}Cotton\textquotedblright{}.) Say it once more.
\end{tcolorbox}
\section[\texorpdfstring{\sk{कौशेयम्} (kauśeyam)}{कौशेयम् (kauśeyam)}]{\sk{कौशेयम्} (kauśeyam) — silk}

\label{lesson:SA-C221-kauseyam}



\begin{quote}
Before the new one: say the Sanskrit for an archer's bow, then the Sanskrit for cotton.
\end{quote}

\subsection*{You'll want to know: \sk{कौशेयम्}}
\textbf{\sk{कौशेयम्}} — \emph{kauśeyam} — \textquotedblleft{}silk\textquotedblright{}.

\begin{grammarlens}[title={What you've built}]
One more word for this chapter.
\end{grammarlens}

\subsection*{Your turn}
\begin{itemize}
\raggedright
  \item \emph{Say it:} \emph{kauśeyam}
  \item \emph{Say it:} \emph{kauśeyam}, once more
  \item \emph{Say it:} say \emph{kauśeyam}
  \item \emph{Recall:} say the Sanskrit for an archer's bow, then the Sanskrit for cotton, then say \emph{kauśeyam} again
\end{itemize}

\begin{tcolorbox}[breakable,colback=teal!4,colframe=teal!35!black,title={Before you move on}]
What does \sk{कौशेयम्} mean? (\textquotedblleft{}Silk\textquotedblright{}.) Say it once more.
\end{tcolorbox}
\section[\texorpdfstring{\sk{चर्म} (carma)}{चर्म (carma)}]{\sk{चर्म} (carma) — leather, hide}

\label{lesson:SA-C221-carma}



\begin{quote}
Before the new one: say the Sanskrit for cotton, then the Sanskrit for silk.
\end{quote}

\subsection*{You'll want to know: \sk{चर्म}}
\textbf{\sk{चर्म}} — \emph{carma} — \textquotedblleft{}leather, hide\textquotedblright{}.

\begin{grammarlens}[title={What you've built}]
One more word for this chapter.
\end{grammarlens}

\subsection*{Your turn}
\begin{itemize}
\raggedright
  \item \emph{Say it:} \emph{carma}
  \item \emph{Say it:} \emph{carma}, once more
  \item \emph{Say it:} say \emph{carma}
  \item \emph{Recall:} say the Sanskrit for cotton, then the Sanskrit for silk, then say \emph{carma} again
\end{itemize}

\begin{tcolorbox}[breakable,colback=teal!4,colframe=teal!35!black,title={Before you move on}]
What does \sk{चर्म} mean? (\textquotedblleft{}Leather, hide\textquotedblright{}.) Say it once more.
\end{tcolorbox}
\section[\texorpdfstring{\sk{ग्रन्थिः} (granthiḥ)}{ग्रन्थिः (granthiḥ)}]{\sk{ग्रन्थिः} (granthiḥ) — a knot}

\label{lesson:SA-C221-granthih}



\begin{quote}
Before the new one: say the Sanskrit for silk, then the Sanskrit for leather.
\end{quote}

\subsection*{You'll want to know: \sk{ग्रन्थिः}}
\textbf{\sk{ग्रन्थिः}} — \emph{granthiḥ} — \textquotedblleft{}a knot\textquotedblright{}.

\begin{grammarlens}[title={What you've built}]
One more word for this chapter.
\end{grammarlens}

\subsection*{Your turn}
\begin{itemize}
\raggedright
  \item \emph{Say it:} \emph{granthiḥ}
  \item \emph{Say it:} \emph{granthiḥ}, once more
  \item \emph{Say it:} say \emph{granthiḥ}
  \item \emph{Recall:} say the Sanskrit for silk, then the Sanskrit for leather, then say \emph{granthiḥ} again
\end{itemize}

\begin{tcolorbox}[breakable,colback=teal!4,colframe=teal!35!black,title={Before you move on}]
What does \sk{ग्रन्थिः} mean? (\textquotedblleft{}A knot\textquotedblright{}.) Say it once more.
\end{tcolorbox}
\section[\texorpdfstring{\sk{पीठिका} (pīṭhikā)}{पीठिका (pīṭhikā)}]{\sk{पीठिका} (pīṭhikā) — a small stool}

\label{lesson:SA-C221-pithika}



\begin{quote}
Before the new one: say the Sanskrit for leather, then the Sanskrit for a knot.
\end{quote}

\subsection*{You'll want to know: \sk{पीठिका}}
\textbf{\sk{पीठिका}} — \emph{pīṭhikā} — \textquotedblleft{}a small stool\textquotedblright{}.

\begin{grammarlens}[title={What you've built}]
One more word for this chapter.
\end{grammarlens}

\subsection*{Your turn}
\begin{itemize}
\raggedright
  \item \emph{Say it:} \emph{pīṭhikā}
  \item \emph{Say it:} \emph{pīṭhikā}, once more
  \item \emph{Say it:} say \emph{pīṭhikā}
  \item \emph{Recall:} say the Sanskrit for leather, then the Sanskrit for a knot, then say \emph{pīṭhikā} again
\end{itemize}

\begin{tcolorbox}[breakable,colback=teal!4,colframe=teal!35!black,title={Before you move on}]
What does \sk{पीठिका} mean? (\textquotedblleft{}A small stool\textquotedblright{}.) Say it once more.
\end{tcolorbox}
